% !PRINTED 675
\weberpage{675}

%<<CONT
Schar haben gleich viel Elemente, und diese Zahl heißt die
\gesperrt{Dimension der Schar} und die Schar heißt eine \gesperrt{$s$-fache}.
Sind $\mathfrak{A}_1, \mathfrak{A}_2, \ldots, \mathfrak{A}_s$ linear abhängig oder unabhängig, so gilt
dasselbe von $\mathfrak{M}\mathfrak{A}_1, \mathfrak{M}\mathfrak{A}_2, \ldots, \mathfrak{M}\mathfrak{A}_s$.

Man kann in einer Schar der Dimension $s > 1$

\weq{12}{S = (\mathfrak{A}_1, \mathfrak{A}_2, \ldots, \mathfrak{A}_s)}

ein Polygon finden, das einen beliebigen Punkt $\mathfrak{P}$ mindestens
einmal öfter enthält als der Teiler $\mathfrak{M}$ dieser Schar.

Dieser Zweck wird erreicht, wenn man für diesen Punkt die
Entwickelung (2) ansetzt und dann eine der Konstanten $c_1, c_2, \ldots, c_s$
aus der Gleichung

\weq{13}{c_1e_1 + c_2e_2 + \cdots + c_se_s = 0}

bestimmt. Diese Polygone bilden dann eine Schar von der $(s-1)$ten
Dimension, deren Teiler durch $\mathfrak{M}\mathfrak{P}$ teilbar ist.

Dieser Satz läßt sich dahin erweitern:

5. \gesperrt{Die Polygone einer $s$-fachen Schar, die durch
ein $r$-Eck $\mathfrak{R}$ teilbar sind, bilden eine mindestens
$(s-r)$-fache Schar.}

Ist dieser Satz schon für ein $r$-Eck bewiesen, so folgt er
für ein $(r+1)$-Eck $\mathfrak{R}\mathfrak{P}$ aus dem vorhergehenden.
Denn durch
Hinzutreten des Punktes $\mathfrak{P}$ wird entweder die Dimension nicht
verändert, wenn $\mathfrak{P}$ im Teiler der durch $\mathfrak{R}$ reduzierten Schar aufgeht,
oder sie wird um 1 vermindert.

Man kann das Polygon $\mathfrak{R}$ so wählen, daß die Dimension von
$S$ genau auf $(s-r)$ reduziert wird. Zu diesem Zweck hat man
die Punkte von $\mathfrak{R}$ successive so zu wählen, daß jeder folgende
im Teiler der durch die vorangehende reduzierten Schar nicht
% [sic] printed 'durch die vorangehende reduzierten Schar', the reference to the preceding points (article die, plural) requires 'vorangehenden'
enthalten ist. Daraus folgt:

6. \gesperrt{In einer $s$-fachen Schar gibt es mindestens
ein Polygon, das durch ein gegebenes $(s-1)$-Eck
teilbar ist.}

Ist die Klasse $A$, der alle Polygone der Schar $S$ angehören,
von der $m$ten Ordnung, so kann die Dimension $s$ der Schar nicht
größer sein als $m+1$. Denn sonst könnte man nach 6. in
$S$ ein Polygon finden, das durch ein $(m+1)$-Eck teilbar wäre,
was widersinnig ist, da $S$ nur $m$-Ecke enthält. Es hat daher die
Dimension $s$ der Schar bei gegebener Klasse $A$ ein \gesperrt{Maximum}.
Ist dieses Maximum erreicht, so sind alle Polygone der Klasse $A$
%CONT>>
